% CS5228 Group 6: English final-report structure.
% Based on the supplied IEEE conference template.
% This is a writing skeleton, not a report of completed final experiments.
% Replace every prompt and TBD before submission. Keep the course layout.
% Main-text budget (including title, abstract, tables and figures):
% Title/abstract 0.4; Introduction 0.6; Data 0.4; EDA/preprocessing 1.4;
% Methods 1.5; Experimental setup 1.0; Results/discussion 2.0;
% Conclusion 0.3 = 7.6 pages, leaving 0.4 pages of buffer.
% References and the appendix are outside the eight-page main-text limit.
% The outline's rendered length does not predict the completed paper's length.

\documentclass[conference]{IEEEtran}
\usepackage{cite}
\usepackage{amsmath,amssymb,amsfonts}
\usepackage{graphicx}
\usepackage{textcomp}
\usepackage{xcolor}
\usepackage{booktabs}
\usepackage{tabularx}
\usepackage{url}

% Writing prompts are visible during drafting. Replace them with report prose.
% Set \outlinefalse only after adding complete prose, figures and results.
\newif\ifoutline
\outlinetrue
\newcommand{\guide}[1]{%
  \ifoutline\par\noindent{\small\color{blue!65!black}%
  \textit{Writing prompt: #1}}\par\fi}
\newcommand{\pending}{\textit{TBD}}
\newcommand{\figureplaceholder}[2]{%
  \fbox{\parbox[c][#1][c]{0.91\linewidth}{%
  \centering\small\textit{Figure placeholder}\\[4pt]#2}}}
\newcommand{\field}[1]{\texttt{\detokenize{#1}}}

\title{Predicting Singapore HDB Rental Prices under Temporal Distribution Shift}
% Revise the title if the completed experiments support a different emphasis.
\author{%
  \IEEEauthorblockN{Ox Alpha (Group 6)}
  \IEEEauthorblockA{%
    Li Chengxuan (A0352428A), Mao Huanqing (A0257997Y)\\
    Sun Zekun (A0351314R), Zhang Ziqing (A0223145R)\\
    National University of Singapore}
}

\begin{document}
\maketitle

\begin{abstract}
\guide{Write approximately 150--180 words after the final experiments are
complete. State the practical rental-prediction problem, the dataset and
temporal challenge, the analysis and preprocessing approach, the evaluated
methods, and the selected final model. Give its verified evaluation results
and the most important interpretation. Describe completed work in the past
tense. Do not present progress-report scores as final results.}
\end{abstract}

\begin{IEEEkeywords}
HDB rental prediction, temporal validation, feature engineering, regression,
spatial auxiliary data
\end{IEEEkeywords}

\section{Introduction}
\label{sec:introduction}
% Approximately 0.6 pages.

\subsection{Motivation and Problem Context}
\guide{Explain how estimating HDB rents supports rental-budget planning and
pricing decisions. Identify the intended users and why prediction is
challenging when market conditions change. Use verified sources for external
market or policy claims. Kaggle ranking should not be the primary motivation.}

\subsection{Research Questions and Contributions}
\guide{Organise the paper around the following questions. State contributions
that the completed experiments actually support, rather than claiming that a
new algorithm was invented.}
\begin{enumerate}
  \item How do temporal, spatial and property patterns inform preprocessing
  and prediction of rental prices?
  \item What predictive value do auxiliary features add beyond the core
  rental attributes under a controlled comparison?
  \item Which completed model improvements produce the most reliable
  performance across later time periods, and where does the final model fail?
\end{enumerate}

\section{Data and Problem Formulation}
\label{sec:data}
% Approximately 0.4 pages.

\subsection{Prediction Task and Core Data}
\guide{Define supervised regression from rental attributes to
\field{MONTHLY_RENT} in SGD. Summarise the supplied inputs, labelled data and
test data. Distinguish raw sample counts from the counts used by the final
training pipeline. The goal is prediction, not estimation of causal effects.}

\begin{table}[t]
  \caption{Overview of the supplied rental data.}
  \label{tab:data}
  \centering
  \small
  \begin{tabularx}{\columnwidth}{lXX}
    \toprule
    Property & Training data & Test data \\
    \midrule
    Raw records & 150,000 & 50,000 \\
    Rental months & Jan. 2021--Mar. 2025 & Mar. 2025--Jul. 2026 \\
    Supplied predictors & 10 & 10 \\
    Target labels & Available & Unavailable \\
    Final modelling records & \pending & 50,000 \\
    \bottomrule
  \end{tabularx}
\end{table}
\guide{Refer to Table~\ref{tab:data}. Explain that a small part of the test
data overlaps the final training month, while most records are later. Do not
describe every test record as strictly later than all training observations.}

\subsection{Auxiliary Data Sources}
\guide{Introduce HDB block metadata, MRT stations, schools, shopping malls,
COE prices and stock prices. Identify which sources were evaluated and which
were retained in the final model. Investigation or feature construction alone
does not establish predictive value. Cite the actual dataset sources.}

\section{Exploratory Data Analysis and Preprocessing}
\label{sec:eda}
% Approximately 1.4 pages. Connect each decision to evidence.

\subsection{Target Distribution and Data Quality}
\guide{Describe the target distribution and raw completeness. Explain why
constant \field{FURNISHED} and \field{FEE} fields were removed. Discuss the
IQR-flagged extremes and the final outlier policy. Do not equate an unusual
rent with a proven recording error. State whether the final model used raw
rent or a transformed target and support the decision with completed tests.}

\subsection{Temporal Patterns and Distribution Shift}
\guide{Explain the rapid rental growth during 2021--2023 followed by much
smaller changes in 2024 and early 2025. Connect this change to the validation
design and the final representation of time. Label early-2025 statistics as
partial-year statistics. Figure~\ref{fig:time} should establish this finding.}
\begin{figure}[t]
  \centering
  \figureplaceholder{25mm}{Monthly mean and median rent, with validation
  boundaries and a clearly labelled late-period plateau.}
  \caption{Rental-price patterns over the training period. Replace this
  placeholder with a plot and specify the sample population in the caption.}
  \label{fig:time}
\end{figure}

\subsection{Spatial and Property Patterns}
\guide{Summarise rent differences by town, flat type and floor area. Explain
the confounding between time, location and property size. Show only the EDA
evidence that supports a modelling decision. Rent per square metre may be
used descriptively but must not enter predictors because it contains the
target. Put additional town and flat-model plots in the appendix.}

\subsection{Core Cleaning and Feature Construction}
\guide{Describe text and date normalisation, equivalent flat-type labels,
location keys, lease-age handling and floor-area transformations. Explain the
purpose of each retained feature group. The current core pipeline produces
19 model inputs, but update this count and the feature descriptions if the
final feature set changes. Keep \field{ROW_ID} and \field{RENT_DATE} for
alignment and validation, and exclude them from the predictor matrix.}

\begin{table}[t]
  \caption{EDA findings and the corresponding preprocessing decisions.}
  \label{tab:preprocessing}
  \centering
  \small
  \begin{tabularx}{\columnwidth}{p{0.31\columnwidth}X}
    \toprule
    Finding & Decision and justification \\
    \midrule
    Constant fields & Remove fields without predictive variation. \\
    Inconsistent category labels & Normalise case, whitespace and equivalent
    room labels. \\
    Temporal regime change & Describe the final trend and seasonal features,
    supported by evaluation. \\
    Invalid derived lease age & Flag and treat as missing, with fold-specific
    imputation when required. \\
    Repeated observations & State and justify the final duplicate policy. \\
    Extreme targets & State the retained or filtered policy and supporting
    evidence. \\
    \bottomrule
  \end{tabularx}
\end{table}

\subsection{Auxiliary Construction and Record Integrity}
\guide{Describe address-key joins, HDB attributes and spatial distances or
counts. Explain how historical availability was handled for MRT and other
snapshot data. For economic variables, describe aggregation, calendar lags
and information availability. Report verified coverage and missingness for
sources actually used.}
\guide{State the final exact-duplicate and repeated-feature policies. The
current core removes 641 exact repeats and retains 149,359 rows. An earlier
auxiliary file collapsed feature vectors and averaged rents. The corrected
Ridge comparison uses the same individual core rows and original targets.
Report the final policy, preserve comparable evaluation populations, and
explain any weighting if records were aggregated. Do not compare scores
computed against original targets with scores computed against group means.}

\section{Data Mining Methods and Model Development}
\label{sec:methods}
% Approximately 1.5 pages. Reserve substantial space for completed improvements.

\subsection{Reference Predictor and Ridge Baseline}
\guide{Describe the historical town--flat-type median reference and its
fallback rules. Explain why Ridge was selected as a fast baseline for sparse,
correlated predictors. State its feature set and target representation.
Describe how it was applied rather than reproducing textbook derivations.}

\subsection{Candidate Models}
\guide{Describe the alternative models actually evaluated after the progress
report. Possible candidates include boosting models, but do not list a method
as an experiment unless it was run. For each model, explain its suitability,
categorical handling, target representation and principal configuration.
Add verified references for advanced methods. Remove unused candidates.}

\subsection{Completed Model Improvements}
\label{sec:improvements}
\guide{Use this subsection for the improvements completed before submission.
For each change, connect an observed weakness to a specific modification and
its controlled evaluation. Report unsuccessful changes when they explain
why a component was excluded. Do not write this subsection as a plan.}

\subsubsection{Temporal Representation and Calibration}
\guide{Describe any completed changes to linear time trends, piecewise or
capped trends, training windows, recency weights or calibration. Explain how
cutoffs and other choices were selected using historical validation. State
which approach the final model retained. Delete this heading if no such
experiment was completed.}

\subsubsection{Feature Interactions and Auxiliary Information}
\guide{Describe completed tests of time--location or property interactions,
alternative categorical representations, and additional auxiliary feature
groups. Identify retained and rejected features. Distinguish a plausible
explanation from an effect demonstrated by ablation.}

\subsubsection{Tuning and Model Combination}
\guide{Describe completed tuning, early stopping or ensemble experiments.
If an ensemble was used, specify its component models, training protocol and
how its weights were selected without using evaluation labels. Remove this
heading when the final method is a single model with no relevant experiment.}

\subsection{Final Model Specification}
\label{sec:finalmodel}
\guide{Give the selected final model or ensemble, exact feature groups,
target transformation, hyperparameters and training procedure. Explain the
selection criterion. This subsection must describe the implemented final
system, not merely the highest-scoring intermediate run. State the final
refitting procedure and any post-processing used for test predictions.}

\section{Experimental Setup}
\label{sec:experiments}
% Approximately 1.0 pages.

\subsection{Chronological Validation}
\guide{Explain why evaluation must train on earlier rentals and predict later
rentals. Table~\ref{tab:folds} shows the existing expanding windows. Replace
or supplement them if the final protocol changes. State the prediction
horizon and any gap. Explain how final performance was evaluated separately
from tuning, or explicitly disclose reuse of the tuning folds.}

\begin{table}[t]
  \caption{Existing chronological validation windows. Update to match the
  final evaluation protocol.}
  \label{tab:folds}
  \centering
  \small
  \begin{tabularx}{\columnwidth}{lXX}
    \toprule
    Fold & Training through & Validation period \\
    \midrule
    1 & Jun. 2023 & Jul.--Dec. 2023 \\
    2 & Dec. 2023 & Jan.--Jun. 2024 \\
    3 & Jun. 2024 & Jul. 2024--Mar. 2025 \\
    \bottomrule
  \end{tabularx}
\end{table}
\guide{The current validation blocks span 6, 6 and 9 months, while the
maximum test prediction distance is 16 months after the last training month.
Describe any completed longer-horizon evaluation, or acknowledge the horizon
limitation. A fold already used to choose models is not an untouched holdout.}

\subsection{Training-Fold Preprocessing and Leakage Controls}
\guide{Describe how numerical imputation, scaling, categorical vocabularies
and infrequent-category grouping were fitted inside each historical training
fold. Explain any model-specific exceptions. Exclude target-derived inputs.
State how unknown categories and auxiliary missingness were handled.}

\subsection{Hyperparameter Selection and Reproducibility}
\guide{Report the actual search strategy, search space, validation criterion,
random seeds, software versions and relevant resource constraints. The
existing Ridge search uses alpha in \{0.1, 1, 10, 100\}; update this if expanded.
Summarise decisive settings in the main text and put the full grid in the
appendix. Every result must be traceable to a saved experiment configuration.}

\subsection{Metrics and Controlled Comparisons}
\guide{Use RMSE as the primary metric and MAE and R-squared as supporting
diagnostics. Define whether results are unweighted means of fold metrics or
pooled validation metrics. Use identical record IDs and original labels for
feature comparisons. State the tuning protocol for model comparisons.
If uncertainty is reported, specify how temporal dependence was handled.}

\section{Final Results and Discussion}
\label{sec:results}
% Approximately 2.0 pages. Keep decisive evidence in the main text.

\subsection{Overall Model Comparison}
\guide{Fill Table~\ref{tab:mainresults} with the final verified results from
one consistent evaluation protocol. Compare the reference, baseline,
improved model and selected final configuration. Explain why the final model
was selected, including performance consistency and relevant trade-offs.
All values below remain TBD until the final runs are frozen.}

\begin{table*}[t]
  \caption{Final comparison under the common evaluation protocol. Replace
  generic model names and remove rows for experiments not completed.}
  \label{tab:mainresults}
  \centering
  \small
  \begin{tabularx}{\textwidth}{lXrrrr}
    \toprule
    Method & Feature set & RMSE (SGD) & MAE (SGD) & $R^2$ & RMSE reduction (\%) \\
    \midrule
    Historical median & Town and flat type & \pending & \pending & \pending & --- \\
    Ridge baseline & Core & \pending & \pending & \pending & \pending \\
    Ridge baseline & Core + HDB/MRT & \pending & \pending & \pending & \pending \\
    Improved model & Specify evaluated inputs & \pending & \pending & \pending & \pending \\
    Final selected model & Specify final inputs & \pending & \pending & \pending & \pending \\
    \bottomrule
  \end{tabularx}
  \par\smallskip
  \parbox{0.95\textwidth}{\footnotesize Specify the score aggregation and
  evaluation population. Define percentage reduction relative to the Core
  Ridge baseline. If the improved and final models are identical, combine
  those rows. Use a clearly separate column for any external competition score.}
\end{table*}

\subsection{Performance across Time Periods}
\guide{Report per-fold results and explain differences between market
periods. A lower average alone is insufficient to establish temporal
robustness. Figure~\ref{fig:modelcomparison} should compare the principal
methods across the same validation windows.}
\begin{figure}[t]
  \centering
  \figureplaceholder{23mm}{Per-fold RMSE for the baseline and final model.}
  \caption{Performance across chronological validation periods. Replace the
  placeholder with results from the final experiment runs.}
  \label{fig:modelcomparison}
\end{figure}

\subsection{Ablation Studies}
\label{sec:ablation}
\guide{Quantify the contribution of completed improvements by adding or
removing one component while keeping the evaluation population fixed. Select
the most informative feature, temporal and modelling ablations. Table~
\ref{tab:ablation} is a template, not a required list of experiments.}
\begin{table}[t]
  \caption{Controlled ablation of the final model.}
  \label{tab:ablation}
  \centering
  \small
  \begin{tabularx}{\columnwidth}{Xrr}
    \toprule
    Configuration & RMSE (SGD) & $\Delta$RMSE \\
    \midrule
    Full final configuration & \pending & 0.00 \\
    Without auxiliary spatial features & \pending & \pending \\
    Without the selected temporal change & \pending & \pending \\
    Without a tested interaction/component & \pending & \pending \\
    \bottomrule
  \end{tabularx}
  \par\smallskip
  \parbox{0.95\columnwidth}{\footnotesize Define $\Delta$RMSE as the
  ablated configuration's score minus the full configuration's score.
  Remove rows for components not tested.}
\end{table}
\guide{The progress report's approximately 0.1\% HDB/MRT gain is an
intermediate Ridge finding. Re-evaluate its contribution under the final
model and describe tiny differences proportionately. Do not call a difference
statistically significant without an appropriate uncertainty analysis.}

\subsection{Error Analysis and Model Interpretation}
\guide{Analyse the final model's errors by month, flat type, town and rent
range. Report subgroup sample sizes, RMSE or MAE, and signed bias. Define
residual as actual minus predicted. Show whether the final model still
overpredicts the late-period plateau and compare with the Ridge baseline.
Explain poor cases with observed evidence, while distinguishing unobserved
causes and hypotheses. Treat correlated coefficients or feature attributions
as predictive descriptions, not causal effects.}
\begin{figure}[t]
  \centering
  \figureplaceholder{25mm}{Monthly actual and predicted mean rent for the
  baseline and final model on the same latest evaluation period.}
  \caption{Temporal calibration of the baseline and final model. Report
  mean bias and the validation period in the completed caption.}
  \label{fig:calibration}
\end{figure}

\subsection{Final Test Predictions and Submission}
\guide{Identify the final refit, submission configuration and 50,000
predictions aligned to the sample-submission IDs. Report an actual competition
score only if available, specifying public/private split and the selected
submission. Do not infer test accuracy from prediction ranges or compare a
competition score directly with a different internal validation population.}

\subsection{Limitations of the Final Findings}
\guide{Discuss limitations supported by the completed analysis: missing
unit-level attributes, small subgroups, duplicate-record uncertainty,
historical snapshot accuracy, correlated features and evaluation-horizon
coverage. Explain how these constrain the claims about the selected model.
Keep this as an assessment of the final evidence, without a future-outlook
or planned-experiments section.}

\section{Conclusion}
\label{sec:conclusion}
% Approximately 0.3 pages.
\guide{Answer the research questions using the completed final results.
Name the selected model, quantify its performance relative to the baseline,
state which improvements and auxiliary features proved useful, and summarise
the most important error or limitation. End with the practical interpretation
of the findings. Do not end with a list of work the team intends to do.}

% References begin on a separate page so that the main-text boundary is clear.
% Check the compiled main text ends within eight pages after replacing prompts.
\clearpage
\begin{thebibliography}{00}
  % Replace these drafting entries with verified citations and cite them in prose.
  % Delete unused entries. Do not retain placeholders in the final submission.
  \bibitem{data_source}
  \textit{Placeholder: verified citation for the course rental dataset and
  the auxiliary data sources actually used.}
  \bibitem{ridge_reference}
  \textit{Placeholder: verified reference for Ridge regression, if cited.}
  \bibitem{final_model_reference}
  \textit{Placeholder: original paper or authoritative reference for the
  advanced method actually used in the final model.}
\end{thebibliography}

\clearpage
\appendices
\section{Team Workload Breakdown}
\label{app:workload}
\guide{Record actual completed contributions. The names, IDs and team name
come from the supplied progress slides. Confirm them before submission.
Only Sun Zekun's core-preprocessing responsibility is prefilled from the
conversation; complete all other entries from the team's actual workload.}
\begin{table*}[t]
  \caption{Team workload and report contributions.}
  \label{tab:workload}
  \centering
  \small
  \begin{tabularx}{\textwidth}{llXX}
    \toprule
    Member & Student ID & Main project responsibilities & Report contributions \\
    \midrule
    Li Chengxuan & A0352428A & \pending & \pending \\
    Mao Huanqing & A0257997Y & \pending & \pending \\
    Sun Zekun & A0351314R & Core data preprocessing; complete other actual
    responsibilities & Core preprocessing; complete actual writing/review tasks \\
    Zhang Ziqing & A0223145R & \pending & \pending \\
    \bottomrule
  \end{tabularx}
\end{table*}

\section{Supplementary Details}
\label{app:supplement}
\guide{Use sparingly. Include only the full hyperparameter search/configuration
table, a compact feature dictionary, or supplementary EDA/error tables that
help reproduce or elaborate the main results. Do not move the final model
definition, validation design, main results or essential interpretation out
of the main text. The paper must remain comprehensible without this appendix.}

\end{document}
